\section{Problem setup and notation}\label{sec:chap4 notation_setup}

Let $V:=[P]$ be the set of basic geographic/voting units in a U.S. state and let $w_1,w_2,\dots,w_P$ be the population in each unit.
The map of the U.S. state induces a graph $G=(V,E)$, where two vertices are connected if they are geographically adjacent.
A congressional plan is a partition of the set of precincts $V$ into $I$ districts, where $I$ is determined by the number of U.S. House seats the state has after apportionment~\cite{eckman2019apportionment}.
At a high level, a plan is admissible if it satisfies a number of constraints.
For instance, it needs to be balanced (up to a certain population deviation threshold) and satisfy contiguity (i.e. geographically connected).

It is convenient to view a district $v$ as a vector that lives in $\mathcal{V}\subset\{0,1\}^P$, where
\begin{align}
v_p =
    \begin{cases}
    1, & \text{if } p\text{ is in the district},\\
    0, & \text{otherwise.}
    \end{cases}
\end{align}
Let $\mathcal{W}:=\mathcal{V}^I/S_I$ be the space of unordered $I$-tuples from $\mathcal{V}$ and let $\mathcal{X}\subset\mathcal{W}$ be the subset of admissible plans.
A natural way to define a distance $\mathcal{V}\times\mathcal{V}\to\mathbb{R}$ between two districts $v$ and $v'$ is via a capped population-weighted $\ell^1$:
\begin{align}\label{eq:dist_district}
    d(v,v') := \min\Big(d_{\max}, \sum_{p=1}^{P}|v_p -v'_p|\,w_p\Big).
\end{align}
We include a distance upper bound $d_{\max}>0$ to reduce the noise when two districts are disjoint.
A common choice for the cap is 
\begin{align}
d_{\max}=\frac{2}{I}\sum_{p=1}^{P} w_p.
\end{align}
We can extend the distance on $\mathcal{V}$ to a distance function on plans $\mathcal{W}$ by taking the minimum over all matchings:
\begin{align}\label{eq:dist_plan}
    D(w,w')
    := \min_{\pi\in S_I}\sum_{i=1}^I d\big(w_i,w'_{\pi(i)}\big).
\end{align}
One of the central themes in this paper is to systematically identify plans that are similar.
Instead of working with plans directly, we start by identifying similar districts and use that information to reduce the complexity of the plan-level problem. 
In the three redistricting plans shown in Figure \ref{fig:letter_district_example}, we see that the districts in the southwest corner of Connecticut (which includes Fairfield County) are similar to each other, while the rest of the state has drastically different districting patterns.
\begin{figure}[!htbp]
    \centering
    \includegraphics[width=1.0\linewidth]{figures/letter_district_example.png}
    \caption{Letter/cluster example based on redistricting samples from Connecticut. From left to right, we have: 1. Fraction of districts containing each precinct; 2. The letter as the centroid of the cluster; 3-5. Plans that contain districts similar to the letter (highlighted in red).}
    \label{fig:letter_district_example}
\end{figure}
\FloatBarrier
To identify districts similar to those in Figure \ref{fig:letter_district_example}, we form $K$ clusters on the district space and reduce the complexity of plans by identifying each district by its centroid.
Given a cluster of districts $\tilde{\mathcal{V}}\subset \mathcal{V}$, we define the centroid $\tilde v$ to be the point in $\{0,1\}^P$ that minimizes the average uncapped weighted $\ell^1$ distance to all points in $\tilde{\mathcal{V}}$.
Under this centroid objective, letters are naturally zero-one valued even if we relax the optimization over $[0,1]^P$.
However, letters are not guaranteed to be an existing district.
An alternative approach is to restrict the distance optimization problem to $\mathcal{V}$.
In this paper, we call these centroids letters.
Using $K$ clusters, we obtain an alphabet of size $K$ and can use these letters to form words, which are unordered tuples of letters of fixed size $I$.
Note that letters are objects in $\{0,1\}^P$ and words are objects in $\mathcal{W}$.

Even after the reduction in district space, we still have an exponentially large word space.
Observe that the space of words relevant to actual plans in the ensemble is much smaller.
For instance, we are not interested in words with repeated letters, assuming that the underlying alphabet forms a good representation of the districts.
Based on this observation, we further reduce the space of words by keeping only the $L$ nearest candidate words for each observed plan $x\in\mathcal{X}$ under $D$ (computed via a beam search with width $L$ in Section~\ref{sec:chap4 methodology}).
We then work with the union of these candidate sets, which yields a manageable representative word set that is convenient for downstream tasks, e.g. stratified sampling.

Let $\mathcal{S}\subset\mathcal{W}$ be a representative set of words, which forms a set of strata over the space of plans.
We construct a soft assignment of plans to strata using a partition of unity (POU) $\{\phi_s\}_{s\in \mathcal{S}}$ over plans $x\in\mathcal{X}$:
\begin{align}\label{eq:pou}
    \sum_{s\in \mathcal{S}}\phi_s(x)=1,\qquad
    \phi_s(x):=\frac{\psi(x;s)}{\sum_{s'\in \mathcal{S}}\psi(x;s')}.
\end{align}
For each plan $x$, let $\mathcal{N}_{\mathcal{S}}(x)\subseteq\mathcal{S}$ denote the retained word neighborhood for that plan, typically the words in $\mathcal{S}$ that appear among its $L$ nearest beam-search candidates.
Note that $\mathcal{N}_{\mathcal{S}}(x)$ is often smaller than the beam-width due to pruning.
We define $\psi$ by
\begin{align}\label{eq:psi_relative}
    \psi(x;s)=
    \begin{cases}
        \exp\!\big(-T\kappa(D(x,s))), & s\in\mathcal{N}_{\mathcal{S}}(x),\\
        0, & \text{otherwise,}
    \end{cases}
\end{align}
where $T\geq 0$ is a temperature parameter controlling sensitivity to distances and $\kappa$ is an $\mathbb{R}\to\mathbb{R}$ function that maps distances (in units of population) into appropriate ranges.
When $T=0$, the function $\psi$ assigns equal unnormalized weights to the retained words of a given plan $x$.
% internal comment: T=0 is what we use for proof of principles. This avoids complication from choice of kappa.
Let $\pi$ be the target distribution on plans.
We define the stratum weights by
\[
    z_s := \mathbb{E}_{\pi}[\phi_s],
\]
and the flux matrix between strata by
\begin{align}\label{eq:flux}
    F_{s,s'} := \mathbb{E}_{\pi_s}[\phi_{s'}]
    = \frac{1}{z_s}\mathbb{E}_{\pi}[\phi_s\phi_{s'}],
    \qquad \text{where } \frac{d\pi_s}{d\pi}=\frac{1}{z_s}\phi_s.
\end{align}
By construction, $F$ is row-stochastic and $z$ is stationary for $F$ (i.e., a left eigenvector with eigenvalue $1$).
Observe that this setup is consistent with the standard US framework, see, e.g. \eqref{eqn: umbrella fixed point}.


We conclude this section with a table of notation.
\begin{table}[!htbp]
    \centering
    \small
\renewcommand{\arraystretch}{1.1}
\begin{tabular}{ll}
    \hline
    Symbol & Meaning \\
    \hline
    $G=(V,E)$ & Precinct adjacency graph \\
    $P$ & Number of precincts ($|V|$) \\
    $V$ & Precinct vertex set ($V=[P]$) \\
    $E$ & Adjacency edge set \\
    $\mathcal{V}$ & District indicator space \\
    $\mathcal{W}$ & Space of unordered $I$-tuples of $\mathcal{V}$ \\
    $\mathcal{X}$ & Admissible plan space \\
    $x$ & A plan in $\mathcal{X}$ \\
    $w_p$ & Population weight of precinct $p$ \\
    $I$ & Number of districts in a plan \\
    $d_{\max}$ & Cap used in $d$ \\
    $d$ & District-level distance \eqref{eq:dist_district} \\
    $D$ & Plan-level distance \eqref{eq:dist_plan} \\
    $K$ & Number of letters (district clusters) \\
    $l$ & A letter (cluster centroid) \\
    $L$ & Word degree (number of candidates kept) \\
    $\mathcal{S}$ & Representative word set / strata \\
    $s$ & A stratum (word) in $\mathcal{S}$ \\
    $w_i(x)$ & $i$-th nearest word to plan $x$ \\
    $\psi(x;s)$ & Unnormalized kernel \eqref{eq:psi_relative} \\
    $\phi_s(x)$ & POU weight \eqref{eq:pou} \\
    $T$ & Temperature parameter in \eqref{eq:psi_relative} \\
    $\pi$ & Target distribution on plans \\
    $\pi_s$ & Stratum distribution \eqref{eq:flux} \\
    $z_s$ & Stratum weight $\be_{\pi}[\phi_s]$ \\
    $F_{s,s'}$ & Flux matrix \eqref{eq:flux} \\
    \hline
\end{tabular}
\caption{Summary of basic notations.}
\label{tab:notation}
\end{table}

